\documentclass[11pt,aspectratio=169]{beamer}
\usepackage{amsmath}
\usepackage{booktabs}
\usepackage{array}

\setbeamertemplate{navigation symbols}{}
\setbeamertemplate{footline}[frame number]
\setbeamersize{text margin left=8mm,text margin right=8mm}

\title{Functions for Leave-One-Out\\Validation in GLMs}
\author{Guangyu Qiu}
\date{}

\begin{document}

\begin{frame}
  \titlepage
\end{frame}

\begin{frame}{Cross-validation: evaluating predictions on held-out data}
\textbf{What happens when we use the same data twice?}
\[
  \widehat\beta = \operatorname{fit}(D),
  \qquad \text{error} = \operatorname{test}(D,\widehat\beta).
\]
\begin{itemize}
  \item The fitted coefficients adapt to both signal and random variation in $D$.
  \item Testing on that same data usually gives optimistic error: predictions look better than they would for new subjects.
\end{itemize}

\medskip
\textbf{A basic solution: separate training and testing data}

Without an external dataset, randomly split the available subjects 2:1.

\medskip
\renewcommand{\arraystretch}{1.2}
\begin{tabular}{@{}p{0.46\textwidth}p{0.46\textwidth}@{}}
\toprule
\textbf{Training set: about $2/3$} & \textbf{Testing set: about $1/3$} \\
\midrule
Estimate the model coefficients. & Predict outcomes using the training coefficients; calculate MAE and RMSE without refitting. \\
\bottomrule
\end{tabular}

\medskip
\textbf{Target:} prediction error for a new subject from the same population.
\end{frame}

\begin{frame}{Cross-validation approaches}
\renewcommand{\arraystretch}{1.25}
\begin{tabular}{@{}p{0.18\textwidth}p{0.43\textwidth}p{0.31\textwidth}@{}}
\toprule
\textbf{Approach} & \textbf{How it uses limited data} & \textbf{Key feature} \\
\midrule
Holdout split & Fit on a training set, evaluate on a separate test set. & Results depend on the chosen split; less training data can reduce power. \\
\midrule
K-fold CV & Split into $K$ roughly equal folds; train on $K-1$ and test on the remaining fold each round. & Each subject is held out and tested exactly once. \\
\midrule
LOO CV & Use $K=n$: hold out one subject per round. & Keeps $n-1$ subjects for training. \\
\bottomrule
\end{tabular}

\medskip
\small LOO is useful when data are limited.\\
LOO requires one fit per subject and can be computationally expensive.
\end{frame}

\begin{frame}{LOO example: five subjects, five fits}
\renewcommand{\arraystretch}{1.13}
\centering
\begin{tabular}{@{}ccc@{}}
\toprule
\textbf{Round} & \textbf{Fit coefficients using subjects} & \textbf{Held-out subject} \\
\midrule
1 & 2, 3, 4, 5 & 1 \\
2 & 1, 3, 4, 5 & 2 \\
3 & 1, 2, 4, 5 & 3 \\
4 & 1, 2, 3, 5 & 4 \\
5 & 1, 2, 3, 4 & 5 \\
\bottomrule
\end{tabular}

\medskip
\[
  \eta_i = x_i^\mathsf{T}\widehat\beta_i,
  \qquad \widehat y_i = g^{-1}(\eta_i),
  \qquad i = 1,\ldots,5.
\]

\small $x_i$: held-out subject $i$'s predictors, including the intercept.\\
Re-estimate $\widehat\beta_i$ each round, the inverse link gives the response-scale prediction $\widehat y_i$.

\medskip
\raggedright
\begin{itemize}
  \item Train on $n-1$ subjects. The held-out subject does not enter that fit.
  \item Collect $n$ predictions, one for each subject.
\end{itemize}
\textbf{This project:} 50 subjects, 50 independent fits, 49 training subjects per fit.
\end{frame}

\begin{frame}{Summarizing LOO prediction errors}
For each held-out subject:
\[
  \text{error} = \text{predicted} - \text{observed}.
\]

\begin{columns}[T,onlytextwidth]
\begin{column}{0.47\textwidth}
\textbf{Mean absolute error (MAE)}
\[
  \mathrm{MAE} = \frac{1}{n}\sum_{i=1}^{n}|\text{error}_i|
\]
\begin{itemize}
  \item Easy to interpret.
  \item Less affected by extreme errors than RMSE.
  \item Best constant prediction: the sample median.
\end{itemize}
\end{column}
\begin{column}{0.47\textwidth}
\textbf{Root mean square error (RMSE)}
\[
  \mathrm{RMSE} = \sqrt{\frac{1}{n}\sum_{i=1}^{n}\text{error}_i^2}
\]
\begin{itemize}
  \item Gives greater weight to large errors.
  \item Best constant prediction: the sample mean.
\end{itemize}
\end{column}
\end{columns}

\bigskip
\small
Both use the outcome's units. Lower is better on the same evaluation data.\\
Different outcome types are supported; binary predictions are probabilities.
\end{frame}

\end{document}
